\subsection{Limitations of the data}

The data design in this thesis has several limitations. This \pinktodo{section/subsection} states these limitations explicitly, and \pinktodo{Chapter 5, Discussion} discusses their implications on the validity and interpretation of the results. 

\begin{itemize}
    \item \textbf{Histopathology focus.} This corpus is filtered to contain only histopathology papers. Therefore, no claims can be made about the results being representative of biomedical literature as a whole, or other biomedical subfields.
    \item \textbf{English-language only.} All papers that were ingested in the database are written in English. Therefore, this thesis does not evaluate whether the document and knowledge-extraction pipeline results transfer to non-English histopathology literature.
    \item \textbf{PMC Open Access bias.} The corpus is drawn from the PMC Open Access papers, which might not represent the full histopathology literature. Consequently, the figure conventions, table structures, layout patterns, and caption styles reflect the article types in the PMC Open Access dataset. 
    \item \textbf{Silver-label limits.} The knowledge-extraction reference labels are generated by a strong language model, rather than by field experts. Therefore, these labels inherit biases of the silver-label generators. Furthermore, the agreement system between the voters, explained in further detail in \pinktodo{Chapter 3, Section 3.5, Agreement-Based Cascading}, depends on cosine-similarity between embeddings, potentially inheriting biases from the embedding models that were used. 
\end{itemize}

These limits constrain what the strict-F1 numbers in \pinktodo{Chapter 4, Results} mean, and how they can be correctly interpreted.

